(1) 取 \(\mathbf u=\mathbf U e^{i(\mathbf k\cdot\mathbf r-\omega t)}\)，代入得 \(\rho\omega^2\mathbf U=\mu k^2\mathbf U+(\lambda+\mu)\mathbf k(\mathbf k\cdot\mathbf U)\)。当 \(\mathbf U\parallel\mathbf k\)，有纵波 \(\omega_P=c_Pk\)，\(c_P=\sqrt{(\lambda+2\mu)/\rho}\)；当 \(\mathbf U\perp\mathbf k\)，横波 \(\omega_S=c_Sk\)，\(c_S=\sqrt{\mu/\rho}\)。

(2) 在三维空间与非零 \(\mathbf k\) 垂直的子空间维数为2，故两个独立偏振具有同一 \(c_S\)，纵波只有1个偏振。速度均不随 \(|k|\) 变化，理想无耗散均匀体弹性波是非色散行波，但同一方向的纵波、横波传播速率不同。

(3) 使用 \(\mu=E/[2(1+\nu)]\)、\(\lambda=E\nu/[(1+\nu)(1-2\nu)]\)，得 \(\boxed{c_P/c_S=\sqrt{2(1-\nu)/(1-2\nu)}}\)。常见热力学稳定的各向同性弹性体要求 \(E>0,-1<\nu<1/2\)，因此 \(\mu>0\) 且体积模量 \(K=\lambda+2\mu/3>0\)。

(4) 弹性机械能密度 \(\mathcal E=\rho|\mathbf u_t|^2/2+\lambda(\nabla\cdot\mathbf u)^2/2+\mu\varepsilon_{ij}\varepsilon_{ij}\)，其中 \(\varepsilon_{ij}=(\partial_i u_j+\partial_j u_i)/2\)。弹性应力 \(\sigma_{ij}=\lambda\varepsilon_{kk}\delta_{ij}+2\mu\varepsilon_{ij}\)，能流 \(S_i=-\sigma_{ij}u_{j,t}\) 满足 \(\partial_t\mathcal E+\nabla\cdot\mathbf S=0\)。单纵/横平面波周期平均能流沿波矢，大小等于其平均能量乘相应 \(c_P\) 或 \(c_S\)；能量正定与实本征频率排除无源体波指数失稳。
